\documentclass[11pt]{article}
\usepackage[margin=1in]{geometry}
\usepackage{amsmath}
\usepackage{booktabs}

\title{Ultra-high-dimensional drift localization: how much does performance decline with $p/n$?}
\author{FSDS benchmark working notes}
\date{}

\begin{document}
\maketitle
\vspace{-1.5em}

\paragraph{Summary.}
Moving from $p/n=20$ to $p/n=100$ ($p=10{,}000\to50{,}000$, $n=500$ per batch, a constant 10\% of the features carrying signal), the decline in selection AUC is \emph{negligible} for the two methods that detect shifted features from the covariates (FSL-Net, DataFix), \emph{moderate} for the random-forest domain classifier, and \emph{not measurable} for the outcome-based methods in the covariate-plus-concept-shift design, where they are already at chance at $p/n=20$. Under \emph{pure concept drift}, where outcome-based methods are the only informative ones, the decline is \emph{large}. A statement of ``slight decline'' is therefore supported for covariate-shift detection only, and should not be extended to concept drift.

\begin{table}[h]
\centering
\caption{Selection AUC and precision@$k$ ($k$ = number of shifted features, chance $=0.10$) in the messy nonlinear design with concept and covariate shift. Means over the 20-point grid of factor strengths $\rho$, one replicate per point (DataFix at $p/n=100$: three grid points).}
\begin{tabular}{lcccc}
\toprule
 & \multicolumn{2}{c}{AUC} & \multicolumn{2}{c}{precision@$k$} \\
Method & $p/n=20$ & $p/n=100$ & $p/n=20$ & $p/n=100$ \\
\midrule
FSL-Net            & 0.748 & 0.747 & 0.513 & 0.512 \\
DataFix            & 0.743 & 0.741 & 0.513 & 0.509 \\
RF domain classifier & 0.728 & 0.622 & 0.487 & 0.303 \\
SGShift $\ell_1$ (specified) & 0.532 & 0.515 & 0.154 & 0.125 \\
PermuCATE          & 0.506 & 0.500 & 0.107 & 0.103 \\
LOCO-R-risk        & 0.498 & 0.496 & 0.121 & 0.118 \\
\bottomrule
\end{tabular}
\end{table}

\paragraph{Why the decline is small for covariate-shift detection.}
These methods exploit the difference between the covariate distributions of the two batches, which is present for half of the signal features at a per-feature strength that does not depend on $p$. A plausible reason for the stability (not tested here) is that FSL-Net and DataFix apply a feature-wise structure to every feature, so the evidence for a given shifted feature does not depend on how many other features are present; FSL-Net was in fact run in blocks of $5000$ features. The random-forest domain classifier degrades more (AUC $0.73\to0.62$), plausibly because a forest with $\sqrt p$ candidate features per split is diluted when most features are noise; we have not tested this explanation. The AUC of FSL-Net and DataFix is stable across the whole grid of $\rho$ (about $0.72$--$0.77$ at $p/n=20$), which indicates that the near-equality at the two values of $p/n$ is not a single-point coincidence.

\paragraph{Why the outcome-based methods show no decline there.}
At $p/n=20$ they are already near chance (AUC $0.50$--$0.53$), so there is nothing left to lose; their weakness in this design is not an ultra-high-dimensional effect.

\paragraph{Where the decline is large.}
Under pure concept drift, the same machinery degrades steeply. In the dedicated PermuCATE study (10 shifted features, five replicates) the AUC falls from $0.76$ at $p/n=20$ to $0.58$ at $p/n=100$ and $0.52$ at $p/n=500$ with random-forest nuisances, and the oracle falls from $0.86$ to $0.79$ and $0.64$. The ablations in the repository attribute most of the gap between the learned and the oracle versions to the accuracy of the outcome regressions, which is limited by the sample size (out-of-sample $R^2$ of $0.05$--$0.15$ for both RF and MLP; a cross-validated lasso with the correct model class reaches only $0.17$ and $0.09$ at $p/n=20$ and $100$).

\paragraph{Limitations.}
(i)~One data-generating process and, in the main benchmark, one replicate per grid point (the grid points differ in $\rho$ and are not independent replicates); (ii)~DataFix at $p/n=100$ uses three grid points only; (iii)~the explanation for the random-forest domain classifier is a conjecture; (iv)~the near-constancy of FSL-Net and DataFix is shown at two values of $p/n$ and may not persist at $p/n=500$, which was not run for these methods; (v)~rows are i.i.d., whereas real panels have dependence that lowers the effective sample size (see the panel note), so the real-data decline may be steeper.

\end{document}
